\documentclass[12pt]{article}
\usepackage[utf8]{inputenc}
\usepackage[T1]{fontenc}
\usepackage{hyperref}
\title{Shaetlan Language Resources for K-12 Curriculum Development}
\author{Compiled by Claude Sonnet 5.5 (claude-sonnet-5-5) for Andr\'as Kornai}
\date{}
\begin{document}
\maketitle

\section{Introduction}
Shaetlan (Shetland dialect, Shetlandic; ISO 639-3 scz) is the Scots-Norn mixed language of the Shetland Islands. It is described as endangered, with declining intergenerational transmission, though still widely spoken. There is a community-led education effort.

\section{Primary Educational Resources}
\begin{itemize}
\item Shaetlan Language Plan 2022 (UHI Shetland), which names support for school materials and primary-level curriculum content among its goals: \url{https://www.shetland.uhi.ac.uk/t4-media/one-web/uhi-shetland-images-and-documents/documents-/shaetlan/Shaetlan-Language-Plan-2022.pdf}
\item Shaetlan-language version of the same plan: \url{https://www.shetland.uhi.ac.uk/t4-media/one-web/uhi-shetland-images-and-documents/documents-/shaetlan/Shaetlan-Language-Plan-2022-Shaetlan.pdf}
\item Using Shetland dialect in education (Shetland.org): \url{https://www.shetland.org/blog/using-shetland-dialect-in-education}
\end{itemize}
A 2010 source says Shetland Islands Council employed a dialect development officer working with ForWirds, and that ForWirds had developed dialect teaching material for nursery and lower primary. This may be out of date.

\section{Linguistic and Reference Materials}
\begin{itemize}
\item Wikipedia, Shetland dialect: \url{https://en.wikipedia.org/wiki/Shetland_dialect}
\item Da Spaektionary, online crowd-sourced dictionary (I Hear Dee, 2022, nearly 8,000 lemmas): \url{https://iheardee.com/spaektionary}
\item Shaetlan grammar and dictionary page (I Hear Dee): \url{https://iheardee.com/english/shaetlan-grammar-dictionary}
\item John J. Graham, \emph{The Shetland Dictionary}, online intro at Shetland ForWirds: \url{https://www.shetlanddialect.org.uk/john-j-grahams-shetland-dictionary-intro}
\item Christie-Johnston, \emph{Shetland Words} (2024 revised edition, over 4,000 entries): \url{https://shop.shetlandtimes.co.uk/products/shetland-words}
\end{itemize}

\section{Religious and Community Resources}
\begin{itemize}
\item Shetland ForWirds, a volunteer group formed in 2004 to promote written and spoken dialect, including in education: \url{https://www.shetlanddialect.org.uk/about-us}
\end{itemize}

\section{Digital and Audio Resources}
\begin{itemize}
\item Da Spaektionary is built on a corpus of Shetland Museum and Archives oral history recordings: \url{https://shetnews.co.uk/2022/07/30/team-behind-new-online-dictionary-of-shaetlan-hopes-for-lively-community-participation}
\item Launch report: \url{https://www.shetlandtimes.co.uk/?p=1136583}
\end{itemize}

\section{Curriculum Development Context}
Shaetlan has the best base of the four languages: dictionaries, an online lexicon, a language plan that explicitly aims at school materials, and a community group active in schools. Gaps found: no complete graded textbook series and no standard orthography was confirmed in the sources I saw. Primary-level materials should be developed with ForWirds and UHI Shetland.
\end{document}